% =========================================================
% Algebraic Laws of Set Operations
% =========================================================

\subsection{Binary and Indexed Set-Algebra Laws}

% ---------------------------------------------------------
% TOOLKIT
% ---------------------------------------------------------

\begin{lratopiccard}{Set Algebra --- Signature}
\begin{lrasignaturetable}
\LraSignatureRow{Carrier/context}{The power set \(\mathcal P(U)\) of a fixed ambient set \(U\).}
\LraSignatureRow{Constants}{\(\varnothing\) and \(U\).}
\LraSignatureRow{Operations}{Binary operations \(\cup,\cap\) and unary complement \((-)^c\).}
\LraSignatureRow{Laws}{Commutativity, associativity, distributivity, identity, absorption, and involution.}
\LraSignatureRow{Structural echo}{\(\mathcal P(U)\) is the prototype Boolean algebra of subsets.}
\end{lrasignaturetable}
\end{lratopiccard}

\begin{lraproofpatterncard}{Algebraic Laws --- Proof Pattern}
\begin{lraproofpatterntable}
\LraProofPatternRow{Identity between set terms}{Take an arbitrary element and rewrite membership using the operation rules.}
\LraProofPatternRow{Dual law}{Prove one law directly, then use set duality when the hypotheses match the duality principle.}
\LraProofPatternRow{Indexed law}{Translate membership in an indexed union/intersection into an existential/universal statement over the index set.}
\LraProofPatternRow{Final step}{Convert equivalent membership conditions back into set equality by Extensionality.}
\end{lraproofpatterntable}
\end{lraproofpatterncard}

\begin{remark*}[Exposition]
The operations $\cup$, $\cap$, and complement satisfy algebraic laws analogous
to those of logical connectives. Together they endow $\mathcal{P}(U)$ with
the structure of a \emph{Boolean algebra}. These laws justify manipulation
of set expressions in later proofs, particularly those involving equivalence
classes, partitions, and functions.
\end{remark*}

\begin{theorembox}{Theorem (Commutativity of Union and Intersection)}
\begin{theorem}[Commutativity of Union and Intersection]\label{thm:commutativity}
Let $A,B$ be sets. Then
\[
A \cup B = B \cup A
\quad\text{and}\quad
A \cap B = B \cap A.
\]
\hyperref[prf:commutativity]{Go to proof.}
\end{theorem}
\end{theorembox}

\begin{remark*}[Standard quantified statement]
$\forall A, B:\; \forall x\,(x \in A \cup B \leftrightarrow x \in B \cup A)$ and $\forall x\,(x \in A \cap B \leftrightarrow x \in B \cap A)$.
\end{remark*}

\begin{remark*}[Predicate reading]
\[
  \operatorname{SetEqual}(A\cup B,B\cup A)
  \wedge
  \operatorname{SetEqual}(A\cap B,B\cap A).
\]
\end{remark*}

\begin{remark*}[Interpretation]
Union and intersection do not depend on the order of their two arguments.
\end{remark*}

\LeanFormalizes{thm:commutativity}{lra-lean}{LRA.VolumeI.Set.Operations.Laws.Union}{UnionCommutative}{pending}
\LeanFormalizes{thm:commutativity}{lra-lean}{LRA.VolumeI.Set.Operations.Laws.Intersection}{IntersectionCommutative}{pending}

\begin{dependencies}
\begin{itemize}
  \item \hyperref[ax:extensionality]{Axiom of Extensionality}
  \item \hyperref[def:union]{Union}
  \item \hyperref[def:intersection]{Intersection}
\end{itemize}
\end{dependencies}

\begin{theorembox}{Theorem (Associativity of Union and Intersection)}
\begin{theorem}[Associativity of Union and Intersection]\label{thm:associativity}
Let $A,B,C$ be sets. Then
\[
(A \cup B) \cup C = A \cup (B \cup C)
\quad\text{and}\quad
(A \cap B) \cap C = A \cap (B \cap C).
\]
\hyperref[prf:associativity]{Go to proof.}
\end{theorem}
\end{theorembox}

\begin{lracategoricalcard}{Categorical Rendering --- Associativity of Union}
\begin{center}
\begin{tikzpicture}[lra categorical diagram]
  \node[lra cat native object] (A) at (0,2)
    {\LraCatObject{A}{\{x:x\in A\}}};
  \node[lra cat native object] (B) at (3,2)
    {\LraCatObject{B}{\{x:x\in B\}}};
  \node[lra cat native object] (C) at (6,2)
    {\LraCatObject{C}{\{x:x\in C\}}};

  \node[lra cat left object] (AB) at (1.5,0.6)
    {\LraCatObject{A\cup B}{\{x:x\in A\lor x\in B\}}};
  \node[lra cat right object] (BC) at (4.5,0.6)
    {\LraCatObject{B\cup C}{\{x:x\in B\lor x\in C\}}};
  \node[lra cat result object] (ABC) at (3,-0.95)
    {\LraCatObject{A\cup B\cup C}{\{x:x\in A\lor x\in B\lor x\in C\}}};

  \draw[lra cat left arrow] (A) -- node[lra cat label,above left]{group \(A,B\)} (AB);
  \draw[lra cat left arrow] (B) -- (AB);
  \draw[lra cat right arrow] (B) -- (BC);
  \draw[lra cat right arrow] (C) -- node[lra cat label,above right]{group \(B,C\)} (BC);
  \draw[lra cat result arrow] (AB) -- node[lra cat label,left]{add \(C\)} (ABC);
  \draw[lra cat result arrow] (BC) -- node[lra cat label,right]{add \(A\)} (ABC);

  \node[lra cat note] at (3,-1.75)
    {both parenthesizations produce the same membership condition};
\end{tikzpicture}

\end{center}
\[
  (A\cup B)\cup C
  =
  A\cup(B\cup C)
  =
  \{x:x\in A\lor x\in B\lor x\in C\}.
\]
\end{lracategoricalcard}

\begin{remark*}[Standard quantified statement]
$\forall A, B, C:\; \forall x\,\bigl((x \in (A \cup B) \cup C) \leftrightarrow (x \in A \cup (B \cup C))\bigr)$, and analogously for $\cap$.
\end{remark*}

\begin{remark*}[Predicate reading]
\[
  \operatorname{SetEqual}((A\cup B)\cup C,A\cup(B\cup C))
  \wedge
  \operatorname{SetEqual}((A\cap B)\cap C,A\cap(B\cap C)).
\]
\end{remark*}

\begin{remark*}[Interpretation]
Union and intersection can be regrouped without changing the resulting set.
\end{remark*}

\LeanFormalizes{thm:associativity}{lra-lean}{LRA.VolumeI.Set.Operations.Laws.Union}{UnionAssociative}{pending}
\LeanFormalizes{thm:associativity}{lra-lean}{LRA.VolumeI.Set.Operations.Laws.Intersection}{IntersectionAssociative}{pending}

\begin{dependencies}
\begin{itemize}
  \item \hyperref[ax:extensionality]{Axiom of Extensionality}
  \item \hyperref[def:union]{Union}
  \item \hyperref[def:intersection]{Intersection}
\end{itemize}
\end{dependencies}

\begin{theorembox}{Theorem (Distributive Laws)}
\begin{theorem}[Distributive Laws]\label{thm:distributivity}
Let $A,B,C$ be sets. Then
\[
A \cap (B \cup C) = (A \cap B) \cup (A \cap C),
\]
\[
A \cup (B \cap C) = (A \cup B) \cap (A \cup C).
\]
\hyperref[prf:distributivity]{Go to proof.}
\end{theorem}
\end{theorembox}

\begin{remark*}[Standard quantified statement]
$\forall A, B, C:\; \forall x\,\bigl(x \in A \cap (B \cup C) \leftrightarrow x \in (A \cap B) \cup (A \cap C)\bigr)$, and dually for $\cup$ over $\cap$.
\end{remark*}

\begin{remark*}[Predicate reading]
\[
  \operatorname{SetEqual}(A\cap(B\cup C),(A\cap B)\cup(A\cap C))
  \wedge
  \operatorname{SetEqual}(A\cup(B\cap C),(A\cup B)\cap(A\cup C)).
\]
\end{remark*}

\begin{remark*}[Interpretation]
Intersection distributes over union, and union distributes over intersection.
\end{remark*}

\LeanFormalizes{thm:distributivity}{lra-lean}{LRA.VolumeI.Set.Operations.Laws.Duality}{IntersectionDistributesOverUnion}{pending}
\LeanFormalizes{thm:distributivity}{lra-lean}{LRA.VolumeI.Set.Operations.Laws.Duality}{UnionDistributesOverIntersection}{pending}

\begin{dependencies}
\begin{itemize}
  \item \hyperref[ax:extensionality]{Axiom of Extensionality}
  \item \hyperref[def:union]{Union}
  \item \hyperref[def:intersection]{Intersection}
\end{itemize}
\end{dependencies}

\begin{theorembox}{Theorem (Distributive Laws for Indexed Families)}
\begin{theorem}[Distributive Laws for Indexed Families]\label{thm:indexed-distributivity}
Let $A\subseteq U$ and let $\{B_i\}_{i\in I}$ be an indexed family of subsets
of a universe $U$. Then
\[
A\cap\bigcup_{i\in I}B_i
=
\bigcup_{i\in I}(A\cap B_i).
\]
If $I\neq\varnothing$, then
\[
A\cup\bigcap_{i\in I}B_i
=
\bigcap_{i\in I}(A\cup B_i).
\]
\hyperref[prf:indexed-distributivity]{Go to proof.}
\end{theorem}
\end{theorembox}

\begin{remark*}[Standard quantified statement]
\[
  A\cap\bigcup_{i\in I}B_i
  =
  \bigcup_{i\in I}(A\cap B_i),
  \qquad
  I\neq\varnothing\to
  A\cup\bigcap_{i\in I}B_i
  =
  \bigcap_{i\in I}(A\cup B_i).
\]
\end{remark*}

\begin{remark*}[Interpretation]
Union and intersection distribute over indexed families in the same way they
distribute over binary operations.
\end{remark*}

\LeanFormalizes{thm:indexed-distributivity}{lra-lean}{LRA.VolumeI.Set.Operations.Laws.Indexed}{IntersectionDistributesOverIndexedUnion}{pending}
\LeanFormalizes{thm:indexed-distributivity}{lra-lean}{LRA.VolumeI.Set.Operations.Laws.Indexed}{UnionDistributesOverIndexedIntersection}{pending}

\begin{dependencies}
\begin{itemize}
  \item \hyperref[def:indexed-family]{Indexed Family of Sets}
  \item \hyperref[def:indexed-union]{Indexed Union}
  \item \hyperref[def:indexed-intersection]{Indexed Intersection}
  \item \hyperref[def:union]{Union}
  \item \hyperref[def:intersection]{Intersection}
\end{itemize}
\end{dependencies}

\begin{theorembox}{Theorem (Identity and Absorption Laws)}
\begin{theorem}[Identity and Absorption Laws]\label{thm:identity-absorption}
Let $A,B$ be sets and $U$ a universe with $A \subseteq U$. Then
\[
A \cup \varnothing = A,
\qquad
\varnothing \cup A = A,
\]
\[
A \cap U = A,
\qquad
U \cap A = A,
\]
\[
A\cup U=U,
\qquad
U\cup A=U,
\]
\[
A\cap\varnothing=\varnothing,
\qquad
\varnothing\cap A=\varnothing,
\]
\[
A \cup (A \cap B) = A,
\qquad
A \cap (A \cup B) = A.
\]
\hyperref[prf:identity-absorption]{Go to proof.}
\end{theorem}
\end{theorembox}

\begin{remark*}[Standard quantified statement]
$\forall A \subseteq U:\; \forall x\,(x \in A \cup \varnothing \leftrightarrow x \in A)$ and $\forall x\,(x \in A \cap U \leftrightarrow x \in A)$.
The symmetric identity laws and the domination/null laws are stated explicitly:
$\varnothing\cup A=A$, $U\cap A=A$, $A\cup U=U$, $U\cup A=U$,
$A\cap\varnothing=\varnothing$, and $\varnothing\cap A=\varnothing$.
Absorption: $\forall A, B:\; \forall x\,(x \in A \cup (A \cap B) \leftrightarrow x \in A)$, and dually.
\end{remark*}

\begin{remark*}[Predicate reading]
\[
  \operatorname{SetEqual}(A\cup\varnothing,A)
  \wedge
  \operatorname{SetEqual}(\varnothing\cup A,A)
  \wedge
  \operatorname{SetEqual}(A\cap U,A)
  \wedge
  \operatorname{SetEqual}(U\cap A,A)
  \wedge
  \operatorname{SetEqual}(A\cup U,U)
  \wedge
  \operatorname{SetEqual}(U\cup A,U)
  \wedge
  \operatorname{SetEqual}(A\cap\varnothing,\varnothing)
  \wedge
  \operatorname{SetEqual}(\varnothing\cap A,\varnothing)
  \wedge
  \operatorname{SetEqual}(A\cup(A\cap B),A)
  \wedge
  \operatorname{SetEqual}(A\cap(A\cup B),A).
\]
\end{remark*}

\begin{remark*}[Interpretation]
The empty set and universe act as identities for the appropriate operations,
and absorption removes redundant nested occurrences of $A$.
\end{remark*}

\LeanFormalizes{thm:identity-absorption}{lra-lean}{LRA.VolumeI.Set.Operations.Laws.Union}{UnionEmpty}{pending}
\LeanFormalizes{thm:identity-absorption}{lra-lean}{LRA.VolumeI.Set.Operations.Laws.Union}{EmptyUnion}{pending}
\LeanFormalizes{thm:identity-absorption}{lra-lean}{LRA.VolumeI.Set.Operations.Laws.Union}{UnionUniversal}{pending}
\LeanFormalizes{thm:identity-absorption}{lra-lean}{LRA.VolumeI.Set.Operations.Laws.Union}{UniversalUnion}{pending}
\LeanFormalizes{thm:identity-absorption}{lra-lean}{LRA.VolumeI.Set.Operations.Laws.Intersection}{IntersectionUniversal}{pending}
\LeanFormalizes{thm:identity-absorption}{lra-lean}{LRA.VolumeI.Set.Operations.Laws.Intersection}{UniversalIntersection}{pending}
\LeanFormalizes{thm:identity-absorption}{lra-lean}{LRA.VolumeI.Set.Operations.Laws.Intersection}{IntersectionEmpty}{pending}
\LeanFormalizes{thm:identity-absorption}{lra-lean}{LRA.VolumeI.Set.Operations.Laws.Intersection}{EmptyIntersection}{pending}
\LeanFormalizes{thm:identity-absorption}{lra-lean}{LRA.VolumeI.Set.Operations.Laws.Duality}{AbsorptionUnionIntersection}{pending}
\LeanFormalizes{thm:identity-absorption}{lra-lean}{LRA.VolumeI.Set.Operations.Laws.Duality}{AbsorptionIntersectionUnion}{pending}

\begin{dependencies}
\begin{itemize}
  \item \hyperref[ax:extensionality]{Axiom of Extensionality}
  \item \hyperref[def:empty-set]{Empty Set}
  \item \hyperref[def:union]{Union}
  \item \hyperref[def:intersection]{Intersection}
  \item \hyperref[def:subset]{Subset}
\end{itemize}
\end{dependencies}

\begin{theorembox}{Theorem (Idempotency of Union and Intersection)}
\begin{theorem}[Idempotency of Union and Intersection]\label{thm:idempotency}
Let $A$ be a set. Then
\[
A\cup A=A
\qquad\text{and}\qquad
A\cap A=A.
\]
\hyperref[prf:idempotency]{Go to proof.}
\end{theorem}
\end{theorembox}

\begin{remark*}[Standard quantified statement]
\[
  \forall A,\quad A\cup A=A \quad\text{and}\quad A\cap A=A.
\]
\end{remark*}

\begin{remark*}[Predicate reading]
\[
  \operatorname{SetEqual}(A\cup A,A)
  \wedge
  \operatorname{SetEqual}(A\cap A,A).
\]
\end{remark*}

\begin{remark*}[Interpretation]
Repeating the same set in a union or intersection does not change the result.
\end{remark*}

\LeanFormalizes{thm:idempotency}{lra-lean}{LRA.VolumeI.Set.Operations.Laws.Union}{UnionIdempotent}{pending}
\LeanFormalizes{thm:idempotency}{lra-lean}{LRA.VolumeI.Set.Operations.Laws.Intersection}{IntersectionIdempotent}{pending}

\begin{dependencies}
\begin{itemize}
  \item \hyperref[ax:extensionality]{Axiom of Extensionality}
  \item \hyperref[def:union]{Union}
  \item \hyperref[def:intersection]{Intersection}
\end{itemize}
\end{dependencies}

\begin{theorembox}{Theorem (Complement Extreme and Complementarity Laws)}
\begin{theorem}[Complement Extreme and Complementarity Laws]\label{thm:complement-laws}
Let $A\subseteq U$. Then
\[
\varnothing^c=U,
\qquad
U^c=\varnothing,
\]
\[
A\cup A^c=U,
\qquad
A^c\cup A=U,
\]
and
\[
A\cap A^c=\varnothing,
\qquad
A^c\cap A=\varnothing.
\]
\hyperref[prf:complement-laws]{Go to proof.}
\end{theorem}
\end{theorembox}

\begin{remark*}[Standard quantified statement]
\[
  \forall A\subseteq U,\quad
  \varnothing^c=U,\quad U^c=\varnothing,\quad
  A\cup A^c=U,\quad A^c\cup A=U,\quad
  A\cap A^c=\varnothing,\quad A^c\cap A=\varnothing.
\]
\end{remark*}

\begin{remark*}[Predicate reading]
\[
  \operatorname{SetEqual}(\varnothing^c,U)
  \wedge
  \operatorname{SetEqual}(U^c,\varnothing)
  \wedge
  \operatorname{SetEqual}(A\cup A^c,U)
  \wedge
  \operatorname{SetEqual}(A^c\cup A,U)
  \wedge
  \operatorname{SetEqual}(A\cap A^c,\varnothing)
  \wedge
  \operatorname{SetEqual}(A^c\cap A,\varnothing).
\]
\end{remark*}

\begin{remark*}[Interpretation]
A set and its complement exhaust the universe and are disjoint.
\end{remark*}

\LeanFormalizes{thm:complement-laws}{lra-lean}{LRA.VolumeI.Set.Operations.Laws.Complement}{ComplementEmpty}{pending}
\LeanFormalizes{thm:complement-laws}{lra-lean}{LRA.VolumeI.Set.Operations.Laws.Complement}{ComplementUniversal}{pending}
\LeanFormalizes{thm:complement-laws}{lra-lean}{LRA.VolumeI.Set.Operations.Laws.Complement}{UnionComplement}{pending}
\LeanFormalizes{thm:complement-laws}{lra-lean}{LRA.VolumeI.Set.Operations.Laws.Complement}{ComplementUnion}{pending}
\LeanFormalizes{thm:complement-laws}{lra-lean}{LRA.VolumeI.Set.Operations.Laws.Complement}{IntersectionComplement}{pending}
\LeanFormalizes{thm:complement-laws}{lra-lean}{LRA.VolumeI.Set.Operations.Laws.Complement}{ComplementIntersection}{pending}

\begin{dependencies}
\begin{itemize}
  \item \hyperref[ax:extensionality]{Axiom of Extensionality}
  \item \hyperref[def:empty-set]{Empty Set}
  \item \hyperref[def:complement]{Complement}
  \item \hyperref[def:union]{Union}
  \item \hyperref[def:intersection]{Intersection}
  \item \hyperref[def:subset]{Subset}
\end{itemize}
\end{dependencies}

\begin{theorembox}{Theorem (Involution of Complement)}
\begin{theorem}[Involution of Complement]\label{thm:involution}
For any $A \subseteq U$,
\[
(A^c)^c = A.
\]
\hyperref[prf:involution]{Go to proof.}
\end{theorem}
\end{theorembox}

\begin{remark*}[Standard quantified statement]
$\forall A \subseteq U:\; \forall x\,(x \in (A^c)^c \leftrightarrow x \in A)$.
\end{remark*}

\begin{remark*}[Predicate reading]
\[
  \operatorname{SetEqual}((A^c)^c,A).
\]
\end{remark*}

\begin{remark*}[Interpretation]
Taking the complement twice returns the original set.
\end{remark*}

\LeanFormalizes{thm:involution}{lra-lean}{LRA.VolumeI.Set.Operations.Laws.Complement}{DoubleComplement}{pending}

\begin{dependencies}
\begin{itemize}
  \item \hyperref[ax:extensionality]{Axiom of Extensionality}
  \item \hyperref[def:complement]{Complement}
\end{itemize}
\end{dependencies}

\begin{theorembox}{Theorem (Difference Laws)}
\begin{theorem}[Difference Laws]\label{thm:difference-laws}
Let $A,B,C$ be sets. Then
\[
A\setminus B=A\cap B^c,
\qquad
A\setminus\varnothing=A,
\qquad
\varnothing\setminus A=\varnothing,
\qquad
A\setminus A=\varnothing,
\]
\[
A\setminus(B\cup C)=(A\setminus B)\cap(A\setminus C),
\qquad
A\setminus(B\cap C)=(A\setminus B)\cup(A\setminus C),
\]
\[
(A\cup B)\setminus C=(A\setminus C)\cup(B\setminus C),
\qquad
(A\cap B)\setminus C=(A\setminus C)\cap(B\setminus C),
\]
and
\[
A\setminus B\subseteq A,
\qquad
(A\setminus B)\cap B=\varnothing.
\]
\hyperref[prf:difference-laws]{Go to proof.}
\end{theorem}
\end{theorembox}

\begin{remark*}[Standard quantified statement]
\[
\begin{gathered}
A\setminus B=A\cap B^c,\quad
A\setminus\varnothing=A,\quad
\varnothing\setminus A=\varnothing,\quad
A\setminus A=\varnothing,\\
A\setminus(B\cup C)=(A\setminus B)\cap(A\setminus C),\quad
A\setminus(B\cap C)=(A\setminus B)\cup(A\setminus C),\\
(A\cup B)\setminus C=(A\setminus C)\cup(B\setminus C),\quad
(A\cap B)\setminus C=(A\setminus C)\cap(B\setminus C),\\
A\setminus B\subseteq A,\quad
(A\setminus B)\cap B=\varnothing.
\end{gathered}
\]
\end{remark*}

\begin{remark*}[Interpretation]
Set difference is intersection with a complement, so its algebra is inherited
from intersection and complement.
\end{remark*}

\LeanFormalizes{thm:difference-laws}{lra-lean}{LRA.VolumeI.Set.Operations.Laws.Difference}{DifferenceAsIntersectionComplement}{pending}
\LeanFormalizes{thm:difference-laws}{lra-lean}{LRA.VolumeI.Set.Operations.Laws.Difference}{DifferenceEmpty}{pending}
\LeanFormalizes{thm:difference-laws}{lra-lean}{LRA.VolumeI.Set.Operations.Laws.Difference}{EmptyDifference}{pending}
\LeanFormalizes{thm:difference-laws}{lra-lean}{LRA.VolumeI.Set.Operations.Laws.Difference}{DifferenceSelf}{pending}
\LeanFormalizes{thm:difference-laws}{lra-lean}{LRA.VolumeI.Set.Operations.Laws.Difference}{DifferenceUnion}{pending}
\LeanFormalizes{thm:difference-laws}{lra-lean}{LRA.VolumeI.Set.Operations.Laws.Difference}{DifferenceIntersection}{pending}
\LeanFormalizes{thm:difference-laws}{lra-lean}{LRA.VolumeI.Set.Operations.Laws.Difference}{UnionDifferenceDistributes}{pending}
\LeanFormalizes{thm:difference-laws}{lra-lean}{LRA.VolumeI.Set.Operations.Laws.Difference}{IntersectionDifferenceDistributes}{pending}
\LeanFormalizes{thm:difference-laws}{lra-lean}{LRA.VolumeI.Set.Operations.Laws.Difference}{DifferenceSubsetLeft}{pending}
\LeanFormalizes{thm:difference-laws}{lra-lean}{LRA.VolumeI.Set.Operations.Laws.Difference}{DifferenceDisjointRight}{pending}

\begin{dependencies}
\begin{itemize}
  \item \hyperref[ax:extensionality]{Axiom of Extensionality}
  \item \hyperref[def:set-difference]{Set Difference}
  \item \hyperref[def:complement]{Complement}
  \item \hyperref[def:empty-set]{Empty Set}
  \item \hyperref[def:union]{Union}
  \item \hyperref[def:intersection]{Intersection}
  \item \hyperref[def:subset]{Subset}
  \item \hyperref[thm:de-morgan]{De Morgan's Laws}
  \item \hyperref[thm:complement-laws]{Complement Extreme and Complementarity Laws}
\end{itemize}
\end{dependencies}

\begin{theorembox}{Theorem (Symmetric Difference Laws)}
\begin{theorem}[Symmetric Difference Laws]\label{thm:symmetric-difference-laws}
Let $A,B,C$ be sets. Then
\[
A\triangle B=(A\setminus B)\cup(B\setminus A),
\qquad
A\triangle B=(A\cup B)\setminus(A\cap B),
\]
\[
A\triangle B=B\triangle A,
\qquad
(A\triangle B)\triangle C=A\triangle(B\triangle C),
\]
\[
A\triangle\varnothing=A,
\qquad
\varnothing\triangle A=A,
\qquad
A\triangle A=\varnothing,
\]
and
\[
A\triangle B=\varnothing\iff A=B,
\qquad
A\triangle B\subseteq A\cup B.
\]
\hyperref[prf:symmetric-difference-laws]{Go to proof.}
\end{theorem}
\end{theorembox}

\begin{remark*}[Standard quantified statement]
\[
\begin{gathered}
A\triangle B=(A\setminus B)\cup(B\setminus A),\quad
A\triangle B=(A\cup B)\setminus(A\cap B),\\
A\triangle B=B\triangle A,\quad
(A\triangle B)\triangle C=A\triangle(B\triangle C),\\
A\triangle\varnothing=A,\quad
\varnothing\triangle A=A,\quad
A\triangle A=\varnothing,\\
A\triangle B=\varnothing\iff A=B,\quad
A\triangle B\subseteq A\cup B.
\end{gathered}
\]
\end{remark*}

\begin{remark*}[Interpretation]
Symmetric difference records exactly where two sets disagree.
\end{remark*}

\LeanFormalizes{thm:symmetric-difference-laws}{lra-lean}{LRA.VolumeI.Set.Operations.Laws.SymmetricDifference}{SymmetricDifferenceAsUnionDifferences}{pending}
\LeanFormalizes{thm:symmetric-difference-laws}{lra-lean}{LRA.VolumeI.Set.Operations.Laws.SymmetricDifference}{SymmetricDifferenceAsUnionDifferenceIntersection}{pending}
\LeanFormalizes{thm:symmetric-difference-laws}{lra-lean}{LRA.VolumeI.Set.Operations.Laws.SymmetricDifference}{SymmetricDifferenceCommutative}{pending}
\LeanFormalizes{thm:symmetric-difference-laws}{lra-lean}{LRA.VolumeI.Set.Operations.Laws.SymmetricDifference}{SymmetricDifferenceEmpty}{pending}
\LeanFormalizes{thm:symmetric-difference-laws}{lra-lean}{LRA.VolumeI.Set.Operations.Laws.SymmetricDifference}{EmptySymmetricDifference}{pending}
\LeanFormalizes{thm:symmetric-difference-laws}{lra-lean}{LRA.VolumeI.Set.Operations.Laws.SymmetricDifference}{SymmetricDifferenceSelf}{pending}
\LeanFormalizes{thm:symmetric-difference-laws}{lra-lean}{LRA.VolumeI.Set.Operations.Laws.SymmetricDifference}{SymmetricDifferenceAssociative}{pending}
\LeanFormalizes{thm:symmetric-difference-laws}{lra-lean}{LRA.VolumeI.Set.Operations.Laws.SymmetricDifference}{SymmetricDifferenceEqEmptyIff}{pending}
\LeanFormalizes{thm:symmetric-difference-laws}{lra-lean}{LRA.VolumeI.Set.Operations.Laws.SymmetricDifference}{SymmetricDifferenceSubsetUnion}{pending}

\begin{dependencies}
\begin{itemize}
  \item \hyperref[ax:extensionality]{Axiom of Extensionality}
  \item \hyperref[def:sym-diff]{Symmetric Difference}
  \item \hyperref[def:set-difference]{Set Difference}
  \item \hyperref[def:empty-set]{Empty Set}
  \item \hyperref[def:union]{Union}
  \item \hyperref[def:intersection]{Intersection}
  \item \hyperref[def:subset]{Subset}
  \item \hyperref[thm:difference-laws]{Difference Laws}
\end{itemize}
\end{dependencies}
